% ============================================================================
%  SECCIÓN 5.2: CLASIFICACIÓN DE LA INFORMACIÓN
% ============================================================================

\section{Clasificaci\'on de la informaci\'on}

La informaci\'on que maneja la aplicaci\'on \nombreApp{} se clasifica
seg\'un su naturaleza, su ciclo de vida y el mecanismo de almacenamiento
requerido. Esta clasificaci\'on permite determinar qu\'e datos se
almacenan en el servidor, cu\'ales en el dispositivo y cu\'ales se
mantienen \'unicamente en memoria.

\subsection{Clasificaci\'on por naturaleza}

\begin{table}[H]
  \centering
  \caption{Clasificaci\'on de la informaci\'on por naturaleza}
  \contenidoConFuente{%
    \begin{tablaAPA}
      \begin{tabular}{@{}p{2.8cm}p{3.4cm}p{3.0cm}p{5.4cm}@{}}
        \toprule
        \textbf{Categor\'ia} & \textbf{Datos} &
        \textbf{Almacenamiento} & \textbf{Justificaci\'on} \\
        \midrule
        Autenticaci\'on & Token JWT, refresh token,
        datos del usuario & AsyncStorage (local) &
        Deben persistir entre sesiones para evitar
        re-login. \\
        \addlinespace
        Dominio & Escenarios, im\'agenes, deportes,
        eventos & Supabase (servidor) &
        Datos compartidos entre usuarios, requieren
        base de datos centralizada. \\
        \addlinespace
        Interacci\'on del usuario & Favoritos, perfil &
        Supabase (servidor) &
        Datos personales que deben sincronizarse
        entre dispositivos. \\
        \addlinespace
        UI & Filtros, b\'usqueda, estado de carga &
        Memoria (React Query / useState) &
        Datos temporales de la interfaz, se regeneran
        en cada sesi\'on. \\
        \bottomrule
      \end{tabular}
    \end{tablaAPA}%
  }{Elaboraci\'on propia en base al an\'alisis del c\'odigo fuente.}
\end{table}

\subsection{Clasificaci\'on por ciclo de vida}

La informaci\'on se clasifica tambi\'en por su duraci\'on en el tiempo:

\begin{enumerate}
  \item \textbf{Permanente:} datos que deben sobrevivir al cierre de la
    aplicaci\'on y al reinicio del dispositivo. Incluyen la sesi\'on del
    usuario (\textit{JWT} y \textit{refresh token}), los favoritos y el
    perfil. Estos datos se almacenan en AsyncStorage (local) o en Supabase
    (servidor).

  \item \textbf{Semi-permanente:} datos que se mantienen mientras la
    aplicaci\'on est\'a activa pero se regeneran al iniciar una nueva
    sesi\'on. Incluyen el cat\'alogo de escenarios, los eventos pr\'oximos
    y las im\'agenes. Se almacenan en el cach\'e de React Query con un
    tiempo de validez de 5 minutos.

  \item \textbf{Temporal:} datos que existen \'unicamente durante una
    interacci\'on espec\'ifica. Incluyen los t\'erminos de b\'usqueda,
    los filtros seleccionados y los estados de carga. Se almacenan en
    el estado local de React (\textit{useState}) y se eliminan al
    cambiar de pantalla.
\end{enumerate}

\subsection{Clasificaci\'on por sensitividad}

Los datos tambi\'en se clasifican por su nivel de sensibilidad:

\begin{table}[H]
  \centering
  \caption{Clasificaci\'on por nivel de sensibilidad}
  \contenidoConFuente{%
    \begin{tablaAPA}
      \begin{tabular}{@{}p{3.0cm}p{5.0cm}p{7.6cm}@{}}
        \toprule
        \textbf{Nivel} & \textbf{Datos} &
        \textbf{Mecanismo de protecci\'on} \\
        \midrule
        Alto & Token JWT, contrase~na, email &
        Cifrado en tr\'ansito (HTTPS), almacenamiento
        seguro en AsyncStorage, tokens de renovaci\'on
        autom\'atica. \\
        \addlinespace
        Medio & Nombre completo, rol, favoritos &
        Acceso restringido por RLS (Row Level Security)
        en Supabase, pol\'iticas por rol. \\
        \addlinespace
        Bajo & Escenarios p\'ublicos, eventos,
        im\'agenes & Datos p\'ublicos, acceso sin
        autenticaci\'on para lectura. \\
        \bottomrule
      \end{tabular}
    \end{tablaAPA}%
  }{Elaboraci\'on propia.}
\end{table}

\subsection{Resumen de la clasificaci\'on}

La siguiente figura resume la clasificaci\'on de la informaci\'on
seg\'un su naturaleza y mecanismo de almacenamiento:

\begin{figure}[H]
\centering
  \figuraAPA{fig:clasificacion-informacion}{Clasificaci\'on de la informaci\'on por almacenamiento}{%
  \begin{tikzpicture}[
    node distance=0.8cm,
    caja/.style={rectangle, draw=black!60, rounded corners,
      minimum width=3.5cm, minimum height=0.7cm, align=center,
      font=\small},
    flecha/.style={-Stealth, thick}
  ]
    \node[caja] (datos) {Datos de la app};
    \node[caja, below left=of datos] (servidor)
      {Servidor (Supabase)};
    \node[caja, below right=of datos] (local)
      {Local (AsyncStorage)};
    \node[caja, below=of servidor] (dominio)
      {Escenarios, eventos};
    \node[caja, below=of local] (auth)
      {Sesi\'on, tokens};
    \draw[flecha] (datos) -- (servidor);
    \draw[flecha] (datos) -- (local);
    \draw[flecha] (servidor) -- (dominio);
    \draw[flecha] (local) -- (auth);
  \end{tikzpicture}
  }{Elaboraci\'on propia.}
\end{figure}
